% ============================================================
%  第一章 三维欧氏空间中的曲线
%  结构沿袭苏步青、胡和生《微分几何》第一章
%  记号约定：位置向量用 \boldsymbol{}，单位向量（T/N/B/e）用
%            \mathbf{}，曲率/挠率等标量不加粗。
% ============================================================

\section{三维欧氏空间与向量函数}

\subsection{三维欧氏空间 $\mathbb{E}^3$}

通常把我们所处的空间称为三维欧氏空间，确切地说，所谓的三维欧氏空间 $\mathbb{E}^3$ 是一个非空的集合，其中的元素称为"点"。它满足以下基本性质：
\begin{enumerate}
	\item 任意两个不同的点唯一地决定了连接它们的直线；
	\item 不在一条直线上的任意三个不同的点唯一地决定了通过这三点的平面；
	\item 在 $\mathbb{E}^3$ 中存在不共面的四个点；
	\item 过直线外任意一点能够作、并且只能作一条直线与已知直线平行。
\end{enumerate}

最后一个断言就是欧氏几何的平行公理。

\subsection{向量运算}

在 $\mathbb{E}^3$ 中，向量运算满足如下法则。设 $\boldsymbol{a}$、$\boldsymbol{b}$、$\boldsymbol{c}$ 为向量，$\lambda$、$\mu$ 为实数。

\textbf{加法运算：}
\begin{equation}
	\boldsymbol{a}+\boldsymbol{b}=\boldsymbol{b}+\boldsymbol{a}
	\label{eq:向量运算:加法交换律}
\end{equation}
\begin{equation}
	(\boldsymbol{a}+\boldsymbol{b})+\boldsymbol{c}=\boldsymbol{a}+(\boldsymbol{b}+\boldsymbol{c})
	\label{eq:向量运算:加法结合律}
\end{equation}

\textbf{数乘运算：}
\begin{equation}
	\lambda(\boldsymbol{a}+\boldsymbol{b})=\lambda\boldsymbol{a}+\lambda\boldsymbol{b}
	\label{eq:向量运算:数乘分配率1}
\end{equation}
\begin{equation}
	(\lambda+\mu)\boldsymbol{a}=\lambda\boldsymbol{a}+\mu\boldsymbol{a}
	\label{eq:向量运算:数乘分配率2}
\end{equation}
\begin{equation}
	\lambda(\mu\boldsymbol{a})=(\lambda\mu)\boldsymbol{a}
	\label{eq:向量运算:数乘结合律}
\end{equation}

\textbf{点乘（数量积）运算：}
\begin{equation}
	\boldsymbol{a}\cdot(\boldsymbol{b}+\boldsymbol{c})=\boldsymbol{a}\cdot\boldsymbol{b}+\boldsymbol{a}\cdot\boldsymbol{c}
	\label{eq:向量运算:点乘分配率}
\end{equation}
\begin{equation}
	\boldsymbol{a}\cdot\boldsymbol{b}=\boldsymbol{b}\cdot\boldsymbol{a}
	\label{eq:向量运算:点乘交换律}
\end{equation}

\textbf{叉乘（向量积）运算：}
\begin{equation}
	\boldsymbol{a}\times(\boldsymbol{b}+\boldsymbol{c})=\boldsymbol{a}\times\boldsymbol{b}+\boldsymbol{a}\times\boldsymbol{c}
	\label{eq:向量运算:叉乘分配率}
\end{equation}
\begin{equation}
	\boldsymbol{a}\times(\boldsymbol{b}\times\boldsymbol{c})=(\boldsymbol{a}\cdot\boldsymbol{c})\boldsymbol{b}-(\boldsymbol{a}\cdot\boldsymbol{b})\boldsymbol{c}
	\label{eq:向量运算:双重向量积}
\end{equation}
\begin{equation}
	\boldsymbol{a}\times\boldsymbol{b}=-(\boldsymbol{b}\times\boldsymbol{a})
	\label{eq:向量运算:叉乘反交换}
\end{equation}
\begin{equation}
	(\lambda\boldsymbol{a})\times\boldsymbol{b}=\boldsymbol{a}\times(\lambda\boldsymbol{b})=\lambda(\boldsymbol{a}\times\boldsymbol{b})
	\label{eq:向量运算:叉乘数乘}
\end{equation}

\textbf{混合积：} 三个向量 $\boldsymbol{a}$、$\boldsymbol{b}$、$\boldsymbol{c}$ 的混合积定义为 $(\boldsymbol{a}\times\boldsymbol{b})\cdot\boldsymbol{c}$。它的几何意义是由 $\boldsymbol{a}$、$\boldsymbol{b}$、$\boldsymbol{c}$ 所张成的平行六面体的\textbf{有向体积}，其结果是一个\textbf{标量}。在不改变三个向量循环次序的前提下，点乘与叉乘可以互换，且三个向量的次序可以进行循环：
\begin{equation}
	(\boldsymbol{a}\times\boldsymbol{b})\cdot\boldsymbol{c}
	=\boldsymbol{a}\cdot(\boldsymbol{b}\times\boldsymbol{c})
	=(\boldsymbol{b}\times\boldsymbol{c})\cdot\boldsymbol{a}
	=\boldsymbol{b}\cdot(\boldsymbol{c}\times\boldsymbol{a})
	\label{eq:向量运算:混合积循环}
\end{equation}

如果交换任意两个相邻向量的顺序，混合积的符号会改变：
\begin{equation}
	(\boldsymbol{a}\times\boldsymbol{b})\cdot\boldsymbol{c}=-(\boldsymbol{b}\times\boldsymbol{a})\cdot\boldsymbol{c}
	\label{eq:向量运算:混合积换序}
\end{equation}

\begin{remark}
	若 $\boldsymbol{a}$、$\boldsymbol{b}$、$\boldsymbol{c}$ 共面，则它们张成的平行六面体体积为零，从而混合积为零。
\end{remark}

\subsection{正交标架}

\begin{definition}[标架与正交标架]
	设 $\mathbf{e}_1$、$\mathbf{e}_2$、$\mathbf{e}_3$ 是 $\mathbb{E}^3$ 中三个不共面的单位向量，定点 $O\in\mathbb{E}^3$，则称 $\{O;\mathbf{e}_1,\mathbf{e}_2,\mathbf{e}_3\}$ 为 $\mathbb{E}^3$ 的一个\textbf{标架}。若 $\mathbf{e}_i$（$i=1,2,3$）两两正交，即
	\begin{equation}
		\mathbf{e}_i\cdot\mathbf{e}_j=\delta_{ij},\qquad i,j=1,2,3,
		\label{eq:正交标架:正交归一条件}
	\end{equation}
	则称该标架为\textbf{右手单位正交标架}，简称为\textbf{正交标架}。
\end{definition}

\begin{theorem}[刚体运动与正交标架]
	$\mathbb{E}^3$ 中的刚体运动把一个正交标架变成一个正交标架；反过来，对于 $\mathbb{E}^3$ 中的任意两个正交标架，必有 $\mathbb{E}^3$ 的一个刚体运动把其中一个正交标架变成另一个正交标架。
\end{theorem}

\subsection{向量函数}

设 $\boldsymbol{u}(t)$、$\boldsymbol{v}(t)$ 是关于参数 $t$ 的可微向量函数，$\phi(t)$ 是一个可微标量函数，$c$ 是一个常数向量或标量，则向量函数的导数满足如下运算法则。

\begin{theorem}[向量函数的求导法则]
	向量函数导数的运算法则如下：
	\begin{enumerate}
		\item 和差的导数（线性性）：
		\begin{equation}
			\frac{\mathrm{d}}{\mathrm{d}t}\bigl[\boldsymbol{u}(t)\pm\boldsymbol{v}(t)\bigr]
			=\boldsymbol{u}'(t)\pm\boldsymbol{v}'(t)
			\label{eq:向量函数:和差导数}
		\end{equation}
		\item 标量函数与向量函数乘积的导数（乘积法则）：
		\begin{equation}
			\frac{\mathrm{d}}{\mathrm{d}t}\bigl[\phi(t)\boldsymbol{u}(t)\bigr]
			=\phi'(t)\boldsymbol{u}(t)+\phi(t)\boldsymbol{u}'(t)
			\label{eq:向量函数:数乘导数}
		\end{equation}
		\item 点积的导数（乘积法则）：
		\begin{equation}
			\frac{\mathrm{d}}{\mathrm{d}t}\bigl[\boldsymbol{u}(t)\cdot\boldsymbol{v}(t)\bigr]
			=\boldsymbol{u}'(t)\cdot\boldsymbol{v}(t)+\boldsymbol{u}(t)\cdot\boldsymbol{v}'(t)
			\label{eq:向量函数:点积导数}
		\end{equation}
		\item 叉积的导数（乘积法则，\textbf{顺序必须保持不变}，因叉积不满足交换律）：
		\begin{equation}
			\frac{\mathrm{d}}{\mathrm{d}t}\bigl[\boldsymbol{u}(t)\times\boldsymbol{v}(t)\bigr]
			=\boldsymbol{u}'(t)\times\boldsymbol{v}(t)+\boldsymbol{u}(t)\times\boldsymbol{v}'(t)
			\label{eq:向量函数:叉积导数}
		\end{equation}
		\item 复合函数的导数（链式法则）：
		\begin{equation}
			\frac{\mathrm{d}}{\mathrm{d}t}\boldsymbol{u}\bigl(s(t)\bigr)
			=\frac{\mathrm{d}\boldsymbol{u}}{\mathrm{d}s}\,\frac{\mathrm{d}s}{\mathrm{d}t}
			=\boldsymbol{u}'\bigl(s(t)\bigr)s'(t)
			\label{eq:向量函数:链式法则}
		\end{equation}
		\item 混合积的导数（展开为三项之和）：
		\begin{equation}
			\frac{\mathrm{d}}{\mathrm{d}t}\bigl[\boldsymbol{u}\cdot(\boldsymbol{v}\times\boldsymbol{w})\bigr]
			=\boldsymbol{u}'\cdot(\boldsymbol{v}\times\boldsymbol{w})
			+\boldsymbol{u}\cdot(\boldsymbol{v}'\times\boldsymbol{w})
			+\boldsymbol{u}\cdot(\boldsymbol{v}\times\boldsymbol{w}')
			\label{eq:向量函数:混合积导数}
		\end{equation}
	\end{enumerate}
\end{theorem}

\begin{remark}
	由式~\eqref{eq:向量函数:数乘导数} 可知，若 $c$ 为常数向量，则
	$\dfrac{\mathrm{d}}{\mathrm{d}t}\bigl[c\boldsymbol{u}(t)\bigr]=c\boldsymbol{u}'(t)$。
\end{remark}

\section{曲线}

\subsection{曲线与正则参数曲线}

\begin{definition}[曲线]
	设 $I$ 是实数轴上的一个区间，映射
	\begin{equation}
		\boldsymbol{r}:I\longrightarrow\mathbb{E}^3,\qquad t\longmapsto\boldsymbol{r}(t)
		\label{eq:曲线:定义}
	\end{equation}
	是 $C^k$（$k\geqslant1$）的，则称 $\boldsymbol{r}(I)$ 为一条曲线，$\boldsymbol{r}(t)$ 称为它的\textbf{参数表示}。
\end{definition}

\begin{definition}[正则点与正则参数曲线]
	若在参数 $t=t_0$ 处一阶导数（切向量）不为零：
	\begin{equation}
		\boldsymbol{r}'(t_0)\neq\boldsymbol{0},
		\label{eq:正则点:条件}
	\end{equation}
	则称 $t_0$ 对应的点为\textbf{正则点}，在该点处的\textbf{切线完全确定}。如果曲线 $\boldsymbol{r}(t)$ 在定义域 $I$ 上处处可导，且所有点都是正则点（即对任意 $t\in I$ 有 $\boldsymbol{r}'(t)\neq\boldsymbol{0}$），则称该曲线为\textbf{正则参数曲线}。
\end{definition}

\begin{remark}
	正则曲线保证了曲线上每一点都有确定的切线，这是微分几何中研究曲线性质的基础。本章如无特别说明，曲线均指正则参数曲线。
\end{remark}

\subsection{切线与法平面}

\begin{proposition}[切线方程]
	在正则点 $\boldsymbol{r}(t_0)$ 处，切线是过点 $\boldsymbol{r}(t_0)$ 且方向向量为 $\boldsymbol{r}'(t_0)$ 的直线，其方程为
	\begin{equation}
		\boldsymbol{R}(s)=\boldsymbol{r}(t_0)+s\boldsymbol{r}'(t_0),
		\label{eq:切线:参数方程}
	\end{equation}
	其中 $\boldsymbol{R}(s)$ 是切线上的任意一点，$s$ 是切线上引入的参数。
\end{proposition}

若记 $\boldsymbol{r}(t_0)=(x_0,y_0,z_0)$、$\boldsymbol{r}'(t_0)=(x'_0,y'_0,z'_0)$，则切线的参数方程为
\begin{equation}
	\begin{cases}
		X=x_0+s x'_0,\\
		Y=y_0+s y'_0,\\
		Z=z_0+s z'_0.
	\end{cases}
	\label{eq:切线:分量方程}
\end{equation}

\begin{definition}[法平面]
	过点 $\boldsymbol{r}(t_0)$ 且以 $\boldsymbol{r}'(t_0)$ 为法向量的平面称为曲线在该点的\textbf{法平面}，其方程为
	\begin{equation}
		\bigl(\boldsymbol{R}-\boldsymbol{r}(t_0)\bigr)\cdot\boldsymbol{r}'(t_0)=0.
		\label{eq:法平面:方程}
	\end{equation}
\end{definition}

\section{弧长}

\begin{theorem}[弧长公式]
	曲线 $C$ 从参数 $t=a$ 到 $t=b$ 之间的弧长 $L$ 为
	\begin{equation}
		L=\int_{a}^{b}\left|\boldsymbol{r}'(t)\right|\,\mathrm{d}t.
		\label{eq:弧长:公式}
	\end{equation}
\end{theorem}

若 $\boldsymbol{r}(t)=\bigl(x(t),y(t),z(t)\bigr)$，则
\begin{equation}
	\left|\boldsymbol{r}'(t)\right|
	=\sqrt{\left(\frac{\mathrm{d}x}{\mathrm{d}t}\right)^{2}
	+\left(\frac{\mathrm{d}y}{\mathrm{d}t}\right)^{2}
	+\left(\frac{\mathrm{d}z}{\mathrm{d}t}\right)^{2}},
	\label{eq:弧长:切向量模}
\end{equation}

弧长元素 $ds$ 是曲线上的微小弧长，它表示参数变化 $\mathrm{d}t$ 内曲线长度的变化量：
\begin{equation}
	ds=\left|\boldsymbol{r}'(t)\right|\,\mathrm{d}t.
	\label{eq:弧长元素:定义}
\end{equation}

\begin{definition}[弧长参数]
	以某一点 $t_0$ 为起点定义函数
	\begin{equation}
		s(t)=\int_{t_0}^{t}\left|\boldsymbol{r}'(\tau)\right|\,\mathrm{d}\tau,
		\label{eq:弧长参数:定义}
	\end{equation}
	则 $s$ 是 $t$ 的严格单调增函数，可反解出 $t=t(s)$。以 $s$ 作为参数的表示 $\boldsymbol{r}(s)$ 称为曲线的\textbf{弧长参数化}（自然参数化），此时
	\begin{equation}
		\left|\boldsymbol{r}'(s)\right|=1.
		\label{eq:弧长参数:单位切向量}
	\end{equation}
\end{definition}

弧长参数化在后续讨论曲率与挠率时起到简化作用。

\section{曲率}

\subsection{单位切向量}

设 $s$ 为弧长参数，曲线 $C$ 在 $s$ 处的单位切向量为
\begin{equation}
	\mathbf{T}(s)=\boldsymbol{r}'(s)=\frac{\mathrm{d}\boldsymbol{r}}{\mathrm{d}s}.
	\label{eq:单位切向量:定义}
\end{equation}

由于 $s$ 是弧长参数，由式~\eqref{eq:弧长参数:单位切向量} 有 $\left|\mathbf{T}(s)\right|=1$。对恒等式 $\mathbf{T}\cdot\mathbf{T}=1$ 两边关于 $s$ 求导，可得
\begin{equation}
	\mathbf{T}(s)\cdot\mathbf{T}'(s)=0,
	\label{eq:单位切向量:正交性}
\end{equation}
即 $\mathbf{T}'(s)$ 恒垂直于 $\mathbf{T}(s)$。

\subsection{曲率}

\begin{definition}[曲率]
	\textbf{曲率} $\kappa(s)$ 度量曲线偏离直线的程度，定义为单位切向量 $\mathbf{T}$ 对弧长 $s$ 的变化率的模长：
	\begin{equation}
		\kappa(s)=\left|\mathbf{T}'(s)\right|
		=\left|\frac{\mathrm{d}\mathbf{T}}{\mathrm{d}s}\right|.
		\label{eq:曲率:定义}
	\end{equation}
\end{definition}

\begin{proposition}[任意参数下的曲率公式]
	对于以任意参数 $t$ 表示的曲线 $\boldsymbol{r}(t)$，曲率公式为
	\begin{equation}
		\kappa(t)=\frac{\left|\boldsymbol{r}'(t)\times\boldsymbol{r}''(t)\right|}
		{\left|\boldsymbol{r}'(t)\right|^{3}}.
		\label{eq:曲率:任意参数公式}
	\end{equation}
\end{proposition}

\begin{definition}[曲率半径]
	当 $\kappa(s)\neq0$ 时，称 $\rho(s)=1/\kappa(s)$ 为曲线在 $s$ 处的\textbf{曲率半径}。
\end{definition}

\section{挠率与 Frenet 标架}

\subsection{Frenet 标架}

Frenet 标架是一个由三个两两正交的单位向量构成的右手坐标系，它随点沿曲线移动而运动。

\begin{definition}[主法向量与次法向量]
	设 $\kappa(s)\neq 0$。\textbf{主法向量} $\mathbf{N}(s)$ 沿着 $\mathbf{T}'(s)$ 的方向，指向曲线弯曲的内侧：
	\begin{equation}
		\mathbf{N}(s)=\frac{\mathbf{T}'(s)}{\left|\mathbf{T}'(s)\right|}
		=\frac{\mathbf{T}'(s)}{\kappa(s)}.
		\label{eq:主法向量:定义}
	\end{equation}
	\textbf{次法向量（副法向量）} $\mathbf{B}(s)$ 与 $\mathbf{T}(s)$、$\mathbf{N}(s)$ 都垂直，由它们的向量积定义以保证构成右手系：
	\begin{equation}
		\mathbf{B}(s)=\mathbf{T}(s)\times\mathbf{N}(s).
		\label{eq:次法向量:定义}
	\end{equation}
\end{definition}

\begin{definition}[Frenet 标架]
	由单位切向量 $\mathbf{T}$、主法向量 $\mathbf{N}$、次法向量 $\mathbf{B}$ 构成的正交归一基：
	\begin{equation}
		\left\{\mathbf{T}(s),\,\mathbf{N}(s),\,\mathbf{B}(s)\right\}
		\label{eq:Frenet标架:定义}
	\end{equation}
	称为曲线的 \textbf{Frenet 标架}。
\end{definition}

\subsection{挠率}

\begin{definition}[挠率]
	\textbf{挠率} $\tau(s)$ 度量曲线偏离平面的程度（即空间扭曲的程度），定义为次法向量 $\mathbf{B}$ 对弧长 $s$ 的变化率的符号模长：
	\begin{equation}
		\tau(s)=-\mathbf{B}'(s)\cdot\mathbf{N}(s)
		=-\left(\frac{\mathrm{d}\mathbf{B}}{\mathrm{d}s}\right)\cdot\mathbf{N}(s).
		\label{eq:挠率:定义}
	\end{equation}
\end{definition}

\begin{proposition}[任意参数下的挠率公式]
	对于以任意参数 $t$ 表示的曲线 $\boldsymbol{r}(t)$，挠率公式为
	\begin{equation}
		\tau(t)=\frac{\bigl(\boldsymbol{r}'(t)\times\boldsymbol{r}''(t)\bigr)\cdot\boldsymbol{r}'''(t)}
		{\left|\boldsymbol{r}'(t)\times\boldsymbol{r}''(t)\right|^{2}}.
		\label{eq:挠率:任意参数公式}
	\end{equation}
\end{proposition}

\begin{theorem}[平面曲线的判定]
	设曲线 $C$ 不是直线，则它是平面曲线当且仅当它的挠率为零：
	\begin{equation}
		\tau(s)=0\iff \text{曲线 }C\text{ 是平面曲线}.
		\label{eq:挠率:平面曲线判定}
	\end{equation}
\end{theorem}

\section{Serret--Frenet 公式}

Serret--Frenet 公式（简称 Frenet 公式）描述了 Frenet 标架三个单位向量沿曲线的变化率，是微分几何的核心公式之一。

\begin{theorem}[Frenet 公式]
	设 $\kappa$ 为曲率，$\tau$ 为挠率，则 Frenet 标架沿曲线的变化满足
	\begin{equation}
		\begin{cases}
			\mathbf{T}'(s)=\kappa\mathbf{N},\\
			\mathbf{N}'(s)=-\kappa\mathbf{T}+\tau\mathbf{B},\\
			\mathbf{B}'(s)=-\tau\mathbf{N},
		\end{cases}
		\label{eq:Frenet公式:方程组}
	\end{equation}
	即矩阵形式
	\begin{equation}
		\begin{pmatrix}
			\mathbf{T}'\\
			\mathbf{N}'\\
			\mathbf{B}'
		\end{pmatrix}
		=
		\begin{pmatrix}
			0 & \kappa & 0\\
			-\kappa & 0 & \tau\\
			0 & -\tau & 0
		\end{pmatrix}
		\begin{pmatrix}
			\mathbf{T}\\
			\mathbf{N}\\
			\mathbf{B}
		\end{pmatrix}.
		\label{eq:Frenet公式:矩阵形式}
	\end{equation}
\end{theorem}

\begin{remark}
	Frenet 公式的系数矩阵是\textbf{反对称矩阵}，这保证了标架 $\{\mathbf{T},\mathbf{N},\mathbf{B}\}$ 在运动过程中始终保持正交归一。
\end{remark}

\section{曲线在一点邻近的形态}

设曲线以弧长为参数，在 $s=s_0$ 处取 Frenet 标架，将 $\boldsymbol{r}(s)$ 在 $s_0$ 处作 Taylor 展开。由式~\eqref{eq:Frenet公式:方程组} 得
\begin{equation}
	\boldsymbol{r}''(s_0)=\kappa\mathbf{N},\qquad
	\boldsymbol{r}'''(s_0)=-\kappa^{2}\mathbf{T}+\kappa'\mathbf{N}+\kappa\tau\mathbf{B},
	\label{eq:局部标准形:导数}
\end{equation}
从而（记 $\Delta s=s-s_0$，在 $s_0$ 处取值）
\begin{equation}
	\boldsymbol{r}(s)-\boldsymbol{r}(s_0)
	=\Delta s\,\mathbf{T}
	+\frac{1}{2}\kappa(\Delta s)^{2}\mathbf{N}
	+\frac{1}{6}\bigl(-\kappa^{2}\mathbf{T}+\kappa'\mathbf{N}+\kappa\tau\mathbf{B}\bigr)(\Delta s)^{3}+o\bigl((\Delta s)^{3}\bigr).
	\label{eq:局部标准形:展开}
\end{equation}

\begin{definition}[密切平面、法平面与从切平面]
	过 $\boldsymbol{r}(s_0)$ 的以下三个平面称为曲线在该点的\textbf{基本平面}：
	\begin{enumerate}
		\item \textbf{密切平面}：由 $\mathbf{T},\mathbf{N}$ 张成；
		\item \textbf{法平面}：由 $\mathbf{N},\mathbf{B}$ 张成（即过该点且垂直于切向的平面）；
		\item \textbf{从切平面}：由 $\mathbf{T},\mathbf{B}$ 张成。
	\end{enumerate}
\end{definition}

由式~\eqref{eq:局部标准形:展开} 可知曲线在一点邻近沿三个基本平面的投影近似为：
\begin{enumerate}
	\item 在密切平面上的投影（略去三阶小量）为抛物线 $y\approx\dfrac{\kappa}{2}(\Delta s)^{2}$，说明曲线总是弯曲到主法向量一侧；
	\item 在法平面上的投影为半立方抛物线；
	\item 在从切平面上的投影为立方抛物线，挠率 $\tau$ 决定了曲线对该平面的偏离。
\end{enumerate}

\section{空间曲线的基本定理}

\begin{theorem}[曲线论基本定理]
	设 $\kappa(s)>0$、$\tau(s)$ 是区间 $I$ 上的两个可微函数，则存在一条以 $s$ 为弧长参数、以 $\kappa(s)$、$\tau(s)$ 分别为曲率和挠率的空间曲线；并且任意两条这样的曲线可以通过一个刚体运动互相重合（即在正交变换与平移的意义下\textbf{唯一}）。
\end{theorem}

\begin{remark}
	该定理说明曲线的局部形状完全由曲率 $\kappa$ 与挠率 $\tau$ 两个函数所确定，即 $\kappa(s)$、$\tau(s)$ 是曲线的内在不变量。定理的构造性证明见附录~\ref{appendix:证明}。
\end{remark}
